% =====================================================================
\chapter{Introduction}\label{ch:intro}
% =====================================================================
\section{About this project}
This is the project report for the course \textbf{CS8.501 -- Research in Information Security}. The project is based on one research paper:
\begin{quote}
H.~Tan, W.~Zheng and P.~Vijayakumar, ``Secure and Efficient Authenticated Key Management Scheme for UAV-Assisted Infrastructure-Less IoVs,''
\emph{IEEE Transactions on Intelligent Transportation Systems}, vol.~24, no.~6, pp.~6389--6400, June 2023~\cite{tan2023}.
\end{quote}

\section{Problem statement}
After an earthquake or a flood, roadside units (RSUs) and mobile towers often stop working. Rescue vehicles still need to communicate, and the
communication must be authenticated and confidential. The base paper uses a \textbf{tethered UAV (TUAV)}, a drone powered through a cable by an
emergency vehicle (ECRV), as a flying base station. Battery UAVs extend its coverage, and a trusted authority (TA) is reached over a satellite link.
The TA authenticates a group of UAVs, and the TUAV gives them a group key.

Our problem has two parts: (i) check whether the scheme really gives the security properties it claims, and (ii) where it does not, design
improvements that keep the system working securely in a disaster area.

\section{Summary of our work}
\begin{center}
\begin{tabularx}{\linewidth}{@{}p{3.2cm} X@{}}
\toprule
\textbf{Part} & \textbf{What we did}\\
\midrule
Study of the paper & We studied the base paper in detail: the problem it solves, its system model, the protocol step by step, and its security and
performance claims (Chapters~\ref{ch:bg}--\ref{ch:p1analysis}).\\
Analysis and design & We checked every security claim against the paper's own equations, proved three weaknesses (F1--F3), and designed four
contributions (C1--C4) that fix them and cover situations the paper does not handle (Chapters~\ref{ch:p2overview}--\ref{ch:concl}).\\
\bottomrule
\end{tabularx}
\end{center}

\section{Team and individual contributions}
All three members studied the base paper, took part in the discussions and reviewed the final report. The main responsibilities were:
\begin{center}
\begin{tabularx}{\linewidth}{@{}l l X@{}}
\toprule
\textbf{Member} & \textbf{Roll number} & \textbf{Main responsibilities}\\
\midrule
Swaraj Kumar & 2025202016 & study of the paper's motivation, system model and protocol; overall framework of our design; compiling the report\\
Hritik Ranjan & 2025202031 & mathematical analysis of the three weaknesses (F1, F2, F3); construction of C1 and its security theorems\\
Harsh Raj & 2025202005 & design and security analysis of C2, C3 and C4; the hardening study A1--A7; security comparison and results\\
\bottomrule
\end{tabularx}
\end{center}

\section{Organisation of the report}
Chapter~\ref{ch:bg} gives the mathematical background. Chapter~\ref{ch:p1paper} summarises the paper's motivation, contributions and related work.
Chapter~\ref{ch:base} describes the protocol, and Chapter~\ref{ch:p1analysis} summarises the paper's own security and performance analysis.
Chapter~\ref{ch:p2overview} gives an overview of our results. Chapters~\ref{ch:f1}--\ref{ch:f3} prove the three weaknesses,
Chapters~\ref{ch:c1}--\ref{ch:hard} present our contributions, and Chapters~\ref{ch:results}--\ref{ch:concl} compare, discuss and conclude.
